\chapter{Developmental Planning, Nationalism and Citizenship}

\section{Developmental Planning and Rural Development}

\begin{KeyConcept}
\begin{itemize}
\item At independence India faced food scarcity (the Bengal famine; population grew 50\% from 1890--1950 but cultivated land only 1.5\%), low literacy, and discrimination. Nehru chose \textbf{planning} (not free capitalism) for \textbf{self-sufficiency} (fearing Lenin's 'economic imperialism') and to \textbf{double the standard of living} --- creating a \textbf{social democracy} (political + economic democracy; the term from Rosa Luxemburg).
\item The \textbf{Planning Commission (1950)} was an executive (not constitutional) body --- to augment material, capital and human resources, formulate five-year plans, and execute them stage-wise. Plans proposed included the Gandhian plan (Gram Swaraj), the Bombay plan (capitalists), and M. N. Roy's People's Plan (collectivisation).
\item The \textbf{mixed economy} separated public (heavy) and private (light) enterprise --- which \textbf{A. R. Desai} called a \textbf{ploy of the Indian bourgeoisie} (the Bombay plan) to avoid competing with MNCs until they had enough capital.
\item \textbf{Community Development Programme (CDP, 1952)} --- aimed to increase agricultural production and transform village social/economic life; eight categories (agriculture, communications, education, health, training, social welfare, supplementary employment, housing). \textbf{Balwant Rai Mehta} (whose committee's report led to Panchayati Raj) redefined 'community' as \emph{all} villagers working together --- not a sectarian term.
\item \textbf{Critics}: \textbf{A. R. Desai} (bureaucratic, top-down, no elected representatives, overlapping departments, no social-service mentality, benefits captured by landlords --- class polarisation); \textbf{S. C. Dube} (70\% of extension benefits to local elites); \textbf{David Mandelbaum} (a widening rich-poor gulf).
\item \textbf{Rural development programmes} (from the 1960s--70s, slogan 'Garibi Hatao') --- now \textbf{targeted} the backward: poverty-alleviation (IRDP), geography-based (drought/desert/hill), sectoral (health/education), and productivity programmes. \textbf{Robert Chambers} called this 'the programmatisation of India'. \textbf{Vivek Chibber} attributes planning's failure to excessive bureaucracy, the capitalist class's organised counter-efforts, the \textbf{import-substitution} model (creating monopolies), and the demobilisation of labour unions. (Contrast: \textbf{South Korea} succeeded with planning, being export-oriented.)
\item \textbf{Post-1991 (neoliberalism)}: state retreats from market (disinvestment), transparency/accountability, and a \textbf{rights-based approach} --- MGNREGA, RTI, RTE (with legislative backing). The Planning Commission was replaced by \textbf{NITI Aayog}.
\end{itemize}
\end{KeyConcept}

\begin{QuickRevision}
\begin{itemize}
\item Nehruvian planning (self-sufficiency; social democracy); Planning Commission 1950; mixed economy (Desai's 'bourgeois ploy').
\item CDP 1952 (production + transformation; Balwant Rai Mehta $\rightarrow$ Panchayati Raj); critics (Desai, Dube 70\% to elites, Mandelbaum).
\item Rural development (targeted; Garibi Hatao; IRDP; Chambers 'programmatisation'; Chibber --- bureaucracy/import-substitution); post-1991 rights-based (MGNREGA/RTI/RTE; NITI Aayog).
\end{itemize}
\end{QuickRevision}

\section{Indian Nationalism: Derivative Discourse and the Civilizational View}

\begin{KeyConcept}
\begin{itemize}
\item Two views of Indian nationalism: it is a \textbf{derivative discourse} of colonialism (\textbf{Partha Chatterjee} --- Indians claimed superiority in the 'spiritual' private sphere while conceding the 'material' public sphere to the British, in three phases: social reform $\rightarrow$ political movement $\rightarrow$ self-rule), or it is \textbf{civilizational}.
\item \textbf{Tagore}: nationalism is \textbf{anti-Indian} --- it creates non-transcendable identities and releases ethno-religious violence; Indian identities were always porous (Kabir, Nanak, Bulleshah). He proposed the idea of \textbf{'homeland'} (one can belong to India without living in its territory).
\item \textbf{Gandhi}: \textbf{spiritual freedom}, \textbf{Gram Swaraj} (decentralised power to the people --- 'Raj Niti' $\rightarrow$ 'Lok Niti'), and \textbf{inter-nationalism} (anti-imperialist, anti-racist, Asian, non-violent).
\item \textbf{Savarkar} ('Hindutva', 1923, inspired by Mazzini): \textbf{Hindutva $\neq$ Hinduism} --- Hindutva is political/militarised; he masculinised the nation (from 'Matribhoomi' to \textbf{Pitrubhoomi} and \textbf{Punyabhoomi} --- fatherland and holy land). Hindus, Jains and Buddhists have both fatherland and holy land in India; Muslims/Christians have only fatherland. His three ideas --- \textbf{Rashtra (nation), Jati (race), Sanskriti (culture)} --- make an \textbf{ethnic/sectarian} nationalism.
\item \textbf{Ambedkar}: nationalism through the \textbf{annihilation of caste} and \textbf{constitutional morality} (not traditional morality). \textbf{Nehru}: secular, industrial, modern (Western) nationalism.
\end{itemize}
\end{KeyConcept}

\begin{QuickRevision}
\begin{itemize}
\item Derivative discourse (Chatterjee --- 3 phases) vs civilizational view.
\item Tagore (anti-nationalist; homeland; Kabir/Nanak/Bulleshah), Gandhi (spiritual freedom; Gram Swaraj; inter-nationalism), Savarkar (Hindutva $\neq$ Hinduism; Pitrubhoomi/Punyabhoomi; Rashtra/Jati/Sanskriti), Ambedkar (constitutional morality), Nehru (secular/industrial).
\end{itemize}
\end{QuickRevision}

\section{Citizenship and Its Four Discourses}

\begin{KeyConcept}
\begin{itemize}
\item \textbf{Citizenship = the right to have rights} --- including the right to make \textbf{competing demands} (the farmers' protest; Kunal Kamra vs Munawar Faruqui as examples of \textbf{limited citizenship} when citizenship becomes ethnic/descent-based).
\item \textbf{Ornit Shani} identifies \textbf{four citizenship discourses} in India: (1) \textbf{liberal} (the individual holds rights --- fundamental rights); (2) \textbf{republican} (rights in accordance with contribution to the \textbf{common good} --- Nehru's idea); (3) \textbf{ethno-nationalist} (citizenship anchored on descent/religion --- the most exclusive); (4) \textbf{non-statist} (Gandhian --- civil disobedience as a right \emph{and} sacred duty; limited state power; e.g. Vinoba Bhave's \textbf{Bhoodan/Gramdan}, and contemporary Gandhian social movements).
\end{itemize}
\end{KeyConcept}

\begin{QuickRevision}
\begin{itemize}
\item Citizenship = the right to have rights
\item competing demands
\item limited citizenship
\item Ornit Shani
\item four citizenship discourses
\item liberal
\item republican
\item common good
\item ethno-nationalist
\item non-statist
\item Bhoodan/Gramdan
\end{itemize}
\end{QuickRevision}

\section{Elite Formation and Civil Society}

\begin{KeyConcept}
\begin{itemize}
\item \textbf{Elites} (unlike the economically dominant) control \textbf{economic, political and ideological} institutions; they are a closed, self-reproducing group. Indian elites: royal families, political dynasties, upper-caste leaders, film/sports stars, pro-state intellectuals, zamindars, corporate and religious leaders; \textbf{new elites} = YouTubers/influencers, start-up founders, media professionals.
\item \textbf{Krishna Kumar} (educational elites): formed by \textbf{early selection} (private English-medium schools --- remnants of the Raj) and \textbf{mass exams} (JEE, UPSC); education has become skill/status acquisition, not knowledge; elites reproduce themselves (two-thirds of civil servants till 1971 were related to civil servants).
\item \textbf{A. R. Desai}: dominant castes (access to land + political power) used education for mobility $\rightarrow$ new elites. \textbf{Beteille}: reservation created a \textbf{'competition for backwardness'} --- political elites emerged within reserved groups. \textbf{Pradeep Chhibber}: \textbf{local elites} gave legitimation to the 1991 economic reforms (and feature in dynastic politics).
\item \textbf{Civil society} (Neera Chandhoke): weak after 1947 (members joined the state), powerful in the 1970s--80s (anti-Emergency, trade-union/caste movements), then weakened after 1991 as \textbf{NGOs} (professional, funded, non-mobilising) replaced movements.
\end{itemize}
\end{KeyConcept}

\begin{QuickRevision}
\begin{itemize}
\item Citizenship = right to have rights (competing demands; limited citizenship --- Kamra/Faruqui); Ornit Shani's 4 discourses (liberal, republican, ethno-nationalist, non-statist).
\item Elites = control economic+political+ideological; Krishna Kumar (early selection + mass exam), Desai (dominant caste), Beteille (competition for backwardness), Chhibber (local elites); new elites (YouTubers, media).
\item Civil society (Chandhoke): weak 1947 $\rightarrow$ powerful 70s--80s $\rightarrow$ weak post-1991 (NGOs).
\end{itemize}
\end{QuickRevision}
